% Source: source/普通高考/2008/2008广东理.pdf, page 1, question 9.
% Vector flowchart; source PDF has no embedded raster image for this figure.
% Nodes: 开始; 输入 m,n; i=1; a=m×i; n 整除 a?; i=i+1; 输出 a,i; 结束.
% Directed edges: start→input→i=1→a=m×i→test;
% test --否→ i=i+1→return to the stem above a=m×i;
% test --是→ 输出(a,i)→结束. The return and output branches are separate.
\begin{tikzpicture}[
  >=Stealth,
  flow/.style={draw, line width=0.55pt, -{Stealth[length=2.1pt,width=1.8pt]},
    align=center, inner sep=0pt, execute at begin node=\setlength{\parindent}{0pt}},
  box/.style={draw, rectangle, minimum height=0.68cm, inner xsep=0.15cm,
    inner ysep=0.08cm, align=center,
    execute at begin node=\setlength{\parindent}{0pt}},
  terminal/.style={box, rounded corners=0.18cm, minimum width=1.35cm},
  io/.style={box, trapezium, trapezium left angle=72, trapezium right angle=108,
    minimum width=2.25cm},
  decision/.style={box, diamond, aspect=2.7, minimum width=3.9cm,
    minimum height=1.05cm, inner sep=0pt},
  branchlabel/.style={align=center, inner sep=0pt,
    execute at begin node=\setlength{\parindent}{0pt}},
]
  \path[use as bounding box] (-3.65,0) rectangle (4.55,11.25);

  \node[terminal] (start) at (0,10.75) {开始};
  \node[io, minimum width=2.45cm] (input) at (0,9.35) {输入 $m,n$};
  \node[box, minimum width=1.55cm] (init) at (0,7.95) {$i=1$};
  \node[box, minimum width=2.80cm] (multiply) at (0,6.45) {$a=m\times i$};
  \node[decision] (test) at (0,4.75) {$n\text{整除}a?$};
  \node[box, minimum width=2.70cm] (increment) at (2.80,6.45) {$i=i+1$};
  \node[io, minimum width=2.70cm] (output) at (0,2.80) {输出 $a,i$};
  \node[terminal] (finish) at (0,1.30) {结束};
  \draw[flow] (start.south) -- (input.north);
  \draw[flow] (input.south) -- (init.north);
  % The return arrow joins the main stem in its interior; keep the stem as one
  % directed edge so no second arrowhead is placed at that join.
  \draw[flow] (init.south) -- (multiply.north);
  \draw[flow] (multiply.south) -- (test.north);
  \draw[flow] (test.south) -- node[branchlabel,right] {是} (output.north);
  \draw[flow] (output.south) -- (finish.north);

  \draw[flow] (test.east) -- node[branchlabel,above] {否} (2.80,4.75)
    -- (2.80,6.00) -- (increment.south);
  \draw[flow] (increment.north) -- (2.80,7.35) -- (0,7.35);
\end{tikzpicture}
